\documentclass[11pt]{article}

\usepackage[margin=1in]{geometry}
\usepackage{amsmath,amssymb}
\usepackage{booktabs}
\usepackage{array}
\usepackage{longtable}
\usepackage[hidelinks]{hyperref}

\title{Supplementary Methods: Calibrating the MOMISH Maternal Model to Kenyan Counties\\
\large Kisii, Makueni, and Mombasa}
\author{}
\date{}

\newcommand{\p}[1]{\texttt{#1}}

\begin{document}
\maketitle

\section{Purpose and scope}

This supplement describes how the maternal health systems model was calibrated
to three Kenyan counties, Kisii, Makueni, and Mombasa. It is written so that the
same procedure can be applied to other counties.

The calibration follows four principles:
\begin{enumerate}
\item Every target comes from a stated data source available for all counties
(the 2022 Kenya DHS and the Kenya Health Information System, KHIS) or from
published literature.
\item Each parameter is solved from an explicit equation that links it to its
target, rather than by numerical search. Stages whose targets depend on one
another are solved together by repeating these equations until the values stop
changing.
\item The same measurement rule is applied to every recorded diagnosis.
\item Parameters are pooled across counties where county data are too sparse or
too inconsistently recorded to support county values.
\end{enumerate}

Locations are home, Level 2/3 (L2/3), Level 4 (L4), and Level 5 (L5). The
\emph{initial} location is where labour starts; the \emph{final} location is
where delivery occurs after any emergency transfer, which is what DHS and KHIS
record. L5 hospitals receive referrals from outside the county, so L5 is
excluded from all rate targets.

\section{Selection of parameters for calibration}

The model contains three broad classes of parameters. First, quantities
that are directly observed for a county, such as ANC coverage, delivery volume,
and facility-level outcome rates, enter as model inputs or calibration targets.
Second, biological and clinical relationships---including relative risks by
maternal risk or anaemia, diagnostic performance, mortality effects of
complications and care, and intervention efficacy---are informed by published
literature. These relationships are not re-estimated merely to improve county
fit; where uncertainty ranges are available, they are sampled in each
simulation run. Third, calibration is reserved for latent or operational
quantities that are not observed directly but determine whether the model
reproduces the observed baseline care pathway and outcomes.

The calibrated quantities therefore describe four aspects of the county health
system: how women are distributed across delivery locations, how often modeled
clinical events occur, how facilities respond through caesarean section or
transfer, and the overall level of maternal mortality. Table~\ref{tab:scope}
states why each quantity requires calibration and whether it is allowed to vary
by county. Parameters held fixed are shown to make clear which mechanisms were
not used to force agreement with county targets.

\begin{longtable}{>{\raggedright\arraybackslash}p{0.18\textwidth} >{\raggedright\arraybackslash}p{0.29\textwidth} >{\raggedright\arraybackslash}p{0.32\textwidth} >{\raggedright\arraybackslash}p{0.13\textwidth}}
\caption{Rationale and treatment of parameters considered during calibration.}\label{tab:scope}\\
\toprule
Model area & Parameters & Why calibrated or held fixed & Treatment \\
\midrule
\endfirsthead
\toprule
Model area & Parameters & Why calibrated or held fixed & Treatment \\
\midrule
\endhead
Preterm birth & \p{PT\_scale} & The model's gestational-age distribution requires a scale factor to reproduce the externally observed facility preterm rate; preterm status subsequently affects risk recognition, planned CS, and transfer. & County-specific \\
Initial delivery location & \p{home\_noANC}, \p{home\_lowrisk}, \p{baseline\_l23\_preference} & These latent routing probabilities translate observed ANC and final delivery-location distributions into initial locations before transfer. & County-specific \\
Recognized-risk routing & \p{OR\_recognized\_L45} & The routing effect of recognized risk is not directly observed by county; a shared value preserves a common behavioral relationship while allowing the baseline preference to vary. & Shared \\
Planned caesarean section & \p{p\_elec\_CS|highrisk}; \p{p\_elec\_CS|preterm} & These probabilities are not reported directly but determine whether the model reproduces planned CS overall and among preterm births. & County; pooled \\
Complication incidence & \p{p\_aph}, \p{p\_ruptured\_uterus}, \p{p\_eclampsia}, \p{p\_OL\_scale}, \p{pph\_incidence\_multiplier}, \p{sepsis\_incidence\_multiplier} & Baseline event probabilities must reproduce each county's observed total complication burden after applying a common measurement rule and pooled cause mix. & County-specific \\
Caesarean capacity & \p{p\_cs\_capacity} [L2/3, L4] & This operational probability represents the ability to perform an indicated emergency CS, rather than the physical proportion of all births a facility could accommodate. & County-specific \\
Emergency transfer & \p{t\_l23\_l45\_preterm}, \p{t\_l23\_l45\_notsevere}, \p{t\_l23\_l45\_predicted}; \p{p\_l5\_given\_l23\_transfer} & Transfer probabilities are not observed directly and must reproduce the post-transfer L2/3 case mix (preterm share, complication ratio) and the observed mix of transfer reasons; the receiving level is taken directly from referral records. & County; shared destination \\
Mortality level & \p{p\_mat\_death\_facility\_baseline}, \p{RR\_MM\_l23\_vs\_l4} & Cause weights, fixed so that modeled deaths reproduce the global distribution of maternal deaths by cause, set relative risk between complications; an overall level and an L2/3 contrast are needed to match observed deaths. & Pooled \\
Held fixed & \p{p\_comp\_severe\_lowrisk}, \p{MM\_weight\_*}, \p{p\_MM\_others}, \p{weight\_facility\_mat}, \p{sen\_comp\_trad}, \p{spec\_comp\_trad}, \p{t\_l23\_l45\_severe} & These represent biological or clinical relationships, diagnostic performance, or assumptions without sufficient county-specific evidence; allowing them to vary would make county fit less identifiable. & Fixed/shared \\
\bottomrule
\end{longtable}

The stages below are numbered in the order they are solved in each calibration
round. Preterm birth (Stage 1) comes first because it changes recognized risk
and therefore affects routing, planned caesarean section, and transfer
eligibility; it is updated from the previous round's simulation. Routing
(Stage 2) then defines where women start labour, and the transfer probabilities
(Stage 3) are solved by re-simulating until stable. Planned caesarean section
and the complication and caesarean capacity parameters (Stages 4 and 5) are
updated from that final simulation. Mortality (Stage 6) is estimated last
because it depends on the resulting complication and care pathways but does
not feed back into them.

\section{Data required for each county}

\begin{longtable}{>{\raggedright\arraybackslash}p{0.34\textwidth} >{\raggedright\arraybackslash}p{0.58\textwidth}}
\caption{Data sources.}\label{tab:sources}\\
\toprule
Source & Quantities \\
\midrule
\endfirsthead
\toprule
Source & Quantities \\
\midrule
\endhead
2022 Kenya DHS, most recent birth within two years (survey weights, clusters, strata) & County 4+ ANC coverage; home delivery by ANC status; L2/3 share by ANC status (check); national ANC gradient \\
KHIS 2022 deliveries by facility level & L2/3 share of facility births; L5 share of L4/5 births \\
KHIS MOH 711 facility-month indicators, January 2022--June 2025 for every county & L2/3 and internal-L4 complications, preterm births, caesarean sections, and maternal deaths \\
Rescue referral records, Kakamega (November 2021--October 2022) & Reasons for L2/3 transfers; receiving level \\
Waiswa et al. (2020), Kenya and Uganda & Facility preterm rate 0.134; share of L4/5 births arriving by referral 0.122 (check) \\
Sandie et al. (2023), sub-Saharan Africa & 40\% of caesarean sections are planned \\
Ghosh et al. (2022), Kenya and Uganda; Nigerian multicentre study \textbf{[citation to add]} & Caesarean section among preterm births 0.265; planned share 1232/4477 \\
Say et al. (2014) & Direct, non-abortion share of maternal deaths 0.618 \\
Clarke-Deelder et al. (2023) & Suspected PPH in Kenyan referral hospitals 9\% (check) \\
\bottomrule
\end{longtable}

The KHIS facility list must be checked for level coding before use; Makueni
County Referral Hospital is recorded as Level 4 but functions as Level 5 and
was reclassified.

\section{Measurement rules}

\paragraph{KHIS processing.} The indicator extract omits facility-months without
reported values. Rates therefore use all facility-months with deliveries as the
denominator and count missing indicator months as zero. Facility-months with
indicator values but no delivery count are excluded. The same period is used
for every county, so that differences between counties are not differences in
time.

\paragraph{Recorded diagnoses.} Every recorded diagnosis (the six maternal
complications, their total, and preterm birth) is multiplied by the
under-recording factor $s=4.48$: the facility preterm rate of Waiswa et al.
(0.134) divided by the 2020 KHIS facility preterm rate in Kakamega County
(2.99\%), where the factor was first derived. Caesarean sections are recorded procedures and are used unscaled.

\paragraph{Cause mix.} Each county keeps its own total complication rate
(\p{p\_mcr\_l4}); the split by cause uses the delivery-weighted mix of the
three counties:
\begin{equation}
T_{k,c}=\p{p\_mcr\_l4}_c\times
\frac{\sum_{c'}D_{c'}R_{k,c'}}{\sum_{c'}D_{c'}\,\p{p\_mcr\_l4}_{c'}},
\end{equation}
where $R_{k,c}$ is the scaled KHIS rate of cause $k$ in county $c$ and $D_c$ the
county's internal-L4 deliveries. Complication burden and recording quality
cannot be separated from KHIS diagnoses alone, but deaths are recorded far more
completely. Deaths per 100 recorded complications are of the same order across
counties (Kisii 0.97, Makueni 1.62, Mombasa 1.06), so county totals are treated
as real. The cause split is not plausible as a true difference: Kisii records 67\% of
its complications as obstructed labour and 17\% as PPH, against 24--38\% and
37--48\% in the other counties. County cause-specific KHIS rates are reported as
checks.

\paragraph{Mortality.} The model represents direct maternal complications only.
KHIS maternal deaths are therefore multiplied by 0.618, the share of maternal
deaths that are direct and not abortion-related (Say et al. 2014).

\section{Calibration procedure}

\begin{longtable}{>{\raggedright\arraybackslash}p{0.17\textwidth} >{\raggedright\arraybackslash}p{0.37\textwidth} >{\raggedright\arraybackslash}p{0.38\textwidth}}
\caption{Calibration stages.}\label{tab:stages}\\
\toprule
Stage & Target & Method \\
\midrule
\endfirsthead
\toprule
Stage & Target & Method \\
\midrule
\endhead
1 Preterm & Preterm among facility births $=0.134$ & Damped proportional update \\
2 Routing & DHS home share by ANC status; KHIS L2/3 share of facility births & Closed form \\
3 Transfers & KHIS L2/3 preterm share; L2/3:L4 complication ratio; Rescue reason mix and receiving level & Counting equations \\
4 Planned CS & $0.40\times$ KHIS L4 CS rate; planned CS among preterm $=0.0729$ (pooled) & Closed form \\
5 Complications and CS capacity & Six internal-L4 complication rates; KHIS CS rates at L2/3 and L4 & Proportional updates \\
6 Mortality & Pooled KHIS L2/3 and internal-L4 MMR (direct causes) & Fixed-point update \\
\bottomrule
\end{longtable}

Stages 1--5 depend on one another, because their targets are measured after
transfer: preterm prevalence changes routing and transfers, transfers change who
delivers at L4, and planned caesarean section and caesarean capacity change who
is monitored in labour and transferred. They are solved together
(Section~\ref{sec:joint}). Mortality does not feed back, so Stage 6 is solved
once afterwards.

\subsection{Risk recognition}

Routing and planned caesarean section depend on whether an ANC mother is
recognized as high risk. With $\pi$ the high-risk prevalence, $Se$ and $Sp$ the
accuracy of risk screening (\p{sen\_risk\_trad}, \p{spec\_risk\_trad}), and
$g_c$ the probability of being dated preterm or post-term,
\begin{equation}
r_c=1-\bigl[1-\pi Se-(1-\pi)(1-Sp)\bigr](1-g_c).
\end{equation}
$g_c$ depends on \p{PT\_scale} and is recomputed each round ($r_c=0.593$--0.594).

\subsection{Stage 1: Preterm birth}

\p{PT\_scale} multiplies the probability of gestational ages below 37 weeks. The
target is the preterm rate among all facility births after transfer,
$T_{PT}=0.134$, common to the three counties. Kisii's KHIS L4 preterm rate is
half that of the other counties while its L2/3 rate is similar, which suggests
under-recording, so county KHIS rates are not used. Because routing and
transfers respond to preterm prevalence, the update is damped and uses the
previous round's simulation:
\begin{equation}
\p{PT\_scale}\leftarrow\p{PT\_scale}\times
\bigl(T_{PT}/\widehat P(PT\mid\text{facility birth})\bigr)^{0.5}.
\end{equation}

\emph{Assumption:} the facility preterm rate is the same in all three counties.

\subsection{Stage 2: Initial delivery location}

Initial location probabilities are
\begin{align}
\text{No ANC:}\quad & P(H)=h_0,\quad P(L23)=(1-h_0)b_c,\\
\text{ANC, not recognized:}\quad & P(H)=h_U,\quad P(L23)=(1-h_U)b_c,\\
\text{ANC, recognized:}\quad & P(H)=0,\quad P(L23)=\frac{b_c}{b_c+\theta(1-b_c)},
\end{align}
with the remainder at L4/5, where $h_0$ is \p{home\_noANC}, $h_U$ is
\p{home\_lowrisk}, $b_c$ is the baseline L2/3 preference
(\p{baseline\_l23\_preference}), and $\theta$ is the odds ratio for L4/5 among
recognized high-risk women (\p{OR\_recognized\_L45}). A share $\tau_c$ of women
starting at L2/3 is transferred (Stage 3).

With 4+ ANC coverage $a_c$, DHS final home shares $F_{H\mid noANC}$ and
$F_{H\mid ANC}$, and the KHIS L2/3 share of facility births $q_c$,
\begin{align}
h_0&=F_{H\mid noANC},\qquad h_U=\frac{F_{H\mid ANC}}{1-r_c},\qquad
H_c=(1-a_c)h_0+a_cF_{H\mid ANC},\\
I_{L23,c}&=\frac{q_c(1-H_c)}{1-\tau_c},
\end{align}
and $b_c$ is the unique root of
\begin{equation}
(1-a_c)(1-h_0)b_c+a_c\left[(1-r_c)(1-h_U)b_c+r_c\frac{b_c}{b_c+\theta(1-b_c)}\right]=I_{L23,c}.
\end{equation}
The derived quantities are $\p{p\_lb\_l23\_45}=\tau_cI_{L23,c}$ and the
monthly births by location (\p{base\_LB}), with L4/5 births split by the KHIS
L5 share.

$\theta=2$ is shared. The same equations applied to the national DHS sample
give $\theta$ between 1.95 and 2.59 depending on the transfer share; county
samples (32--48 births without ANC) are too small to estimate it.

\emph{Assumptions:} DHS is used for home delivery and ANC and KHIS for the facility
split, because KHIS is a facility census with coded levels and the facility
outcome targets come from KHIS. The DHS L2/3 share among ANC mothers is a check;
it agrees in Kisii and Mombasa but not in Makueni (DHS 36\%, model 16\%).

\subsection{Stage 3: Emergency transfer from L2/3}

A woman who starts labour at L2/3 can be transferred, in order, for a predicted
labour emergency without local caesarean capacity
(\p{t\_l23\_l45\_predicted}), a severe complication (\p{t\_l23\_l45\_severe},
90\%), preterm labour without complication (\p{t\_l23\_l45\_preterm}), or a
non-severe complication (\p{t\_l23\_l45\_notsevere}). Transferred women are
received at L5 with probability \p{p\_l5\_given\_l23\_transfer} $=0.579$ (Rescue,
499 transfers with a recorded destination).

One simulation divides women who start at L2/3 and are not moved by the first
two pathways into $A$ (preterm without complication, $E_{pt}$ women), $B$
(non-severe complication, $E_{ns}$ women, preterm share $\pi_B$, $e_B$ events
each), and $C$ (all others). With $p_{pt}$ and $p_{ns}$ the two transfer
probabilities and $s_4$ the share received at L4,
\begin{align}
\widehat P(PT\mid L23)&=\frac{PT_C+E_{pt}(1-p_{pt})+E_{ns}(1-p_{ns})\pi_B}{N_C+E_{pt}(1-p_{pt})+E_{ns}(1-p_{ns})},\\
\widehat R_{L23:L4}&=\frac{\bigl[ev_C+E_{ns}(1-p_{ns})e_B\bigr]/N_{L23}}
{\bigl[ev_{L4}^{other}+s_4E_{ns}p_{ns}e_B\bigr]/\bigl[N_{L4}^{other}+s_4(E_{pt}p_{pt}+E_{ns}p_{ns})\bigr]},
\end{align}
which are solved for $(p_{pt},p_{ns})$ against the KHIS L2/3 preterm share and
the L2/3:L4 complication ratio. No KHIS indicator identifies
\p{t\_l23\_l45\_predicted}, so it is set so that labour emergencies are the
Rescue share $w=0.583$ of transfers with a known reason:
\begin{equation}
\text{labour-emergency transfers}=\frac{w}{1-w}\,
\bigl(\text{severe}+\text{preterm}+\text{non-severe transfers}\bigr).
\end{equation}
Each pass re-simulates to update the counts.

\emph{Assumptions:} the Rescue reason mix and receiving level, observed in
Kakamega County, apply to all three counties; \p{t\_l23\_l45\_severe} is held
fixed; the L4/5 capacity limit is not applied at baseline, because the capacity input
equals the observed L4/5 delivery load and would otherwise ration emergency
transfers. The share of L4/5 births arriving by
referral (0.122, Waiswa et al.) is a share of L4/5 births, not of all births, and
is used only as a check.

The resulting transfer share $\tau$ (0.35--0.44) is above referral rates
reported from primary facilities (about 10--15\%). It is nevertheless supported
by two independent checks: the share of L4/5 births arriving by transfer
(0.11--0.22, against 0.122 in Waiswa et al.) and the KHIS L2/3 preterm share,
both of which fail when $\tau$ is fixed lower (Section~\ref{sec:sensitivity}).
About 60\% of $\tau$ consists of labour emergencies.

\subsection{Stage 4: Planned caesarean section}

Planned caesarean section is assigned to ANC mothers who start labour at L4/5
and are predicted high risk (\p{p\_elec\_CS|highrisk}, $q_{H,c}$, per county) or
predicted preterm (\p{p\_elec\_CS|preterm}, $q_P$, shared). With $a_c$ and $b_c$
the shares of these two groups among final L4 births, and $c_c$, $d_c$ their
shares among preterm women starting at L4/5 (weights $w_c$),
\begin{align}
a_cq_{H,c}+b_cq_P&=0.40\times\text{KHIS L4 CS rate}_c,\\
\textstyle\sum_cw_c\left(c_cq_{H,c}+d_cq_P\right)&=0.265\times1232/4477=0.0729.
\end{align}
Substituting the first equation into the second leaves one linear equation in
$q_P$.

\emph{Assumptions:} 40\% of caesarean sections are planned (sub-Saharan
average); 0.265 is caesarean section of any type among preterm births and is
converted to planned caesarean section with the Nigerian planned share.

\subsection{Stage 5: Complications and caesarean capacity}

Each incidence parameter mainly drives one target, and the modeled rate is
proportional to it, so each is updated by
$x\leftarrow x\,T_x/\widehat R_x$. Caesarean capacity at L4,
\p{p\_cs\_capacity} [L4], is the probability of emergency caesarean section
among women flagged in labour at L4. It is solved from the L4 caesarean target
after subtracting planned caesarean sections and emergency caesarean sections
among women transferred in:
\begin{equation}
\kappa_{L4}\leftarrow\kappa_{L4}\,
\frac{T_{CS,L4}-\widehat E_{planned}-\widehat E_{in}}{\widehat E_{starters}},
\end{equation}
where $\widehat E_{starters}$ is emergency caesarean section among women who
started at L4/5. \p{p\_cs\_capacity} [L2/3] is updated proportionally to the
L2/3 caesarean rate.

\emph{Assumptions:} every woman transferred for a labour emergency receives
emergency caesarean section; the severity probability
(\p{p\_comp\_severe\_lowrisk}) is not estimated, and facility severe maternal
outcomes (SMO) are used as a check.

\subsection{Stage 6: Maternal mortality}

A woman's facility death probability is set by her most lethal severe
complication:
\begin{equation}
p_{i}=\max\Bigl(\max_k S_{ik}\,\p{MM\_weight}_k,\ \p{p\_MM\_others}\Bigr)
\,W_j\,\p{p\_mat\_death\_facility\_baseline},
\end{equation}
where $S_{ik}$ indicates severe complication $k$, and
$W_j=\p{weight\_facility\_mat}\times\p{RR\_MM\_l23\_vs\_l4}$ at L2/3 and
\p{weight\_facility\_mat} at L4 and L5, followed by odds-ratio adjustments for
caesarean section, emergency caesarean section, and transfer delay.
\p{p\_mat\_death\_facility\_baseline} sets the baseline facility risk; skilled-worker
density enters only as the relative change produced by staffing
interventions, because absolute county staffing counts are uncertain.

County deaths are too few to fit county values (1--35 deaths per level), so the
baseline and the L2/3 relative risk are fitted to deaths pooled over the three
counties: 84 deaths in 168,887 internal-L4 deliveries and 7 in 60,657 L2/3
deliveries, which after the direct-cause scaling give 30.8 and 7.1 per 100,000.
Expected deaths are proportional to the baseline, and L2/3 deaths to the
relative risk, so each iteration applies
\begin{align}
A&=T_{L4}/\widehat M_{L4},\qquad
\p{p\_mat\_death\_facility\_baseline}\leftarrow\p{p\_mat\_death\_facility\_baseline}\times A,\\
\p{RR\_MM\_l23\_vs\_l4}&\leftarrow\p{RR\_MM\_l23\_vs\_l4}\times\frac{T_{L23}/\widehat M_{L23}}{A},
\end{align}
where $\widehat M$ is the modeled expected MMR (the sum of individual death
probabilities, which has less Monte Carlo noise than sampled deaths). Three
iterations suffice. Each county is then checked against its own KHIS MMR.

\emph{Assumptions:} cause-specific weights, \p{p\_MM\_others}, and
\p{weight\_facility\_mat} are held fixed; one mortality level applies to
the three counties.

\subsection{Joint solution of Stages 1--5}
\label{sec:joint}

Each round updates \p{PT\_scale} from the previous round's simulation (Stage 1),
applies the Stage 2 closed form with the current $r_c$ and $\tau_c$, and applies
the Stage 3 counting equations until stable; the Stage 4 and 5 updates are then
computed from the same final simulation. Rounds repeat until $\tau$
changes by less than 0.005, \p{PT\_scale} by less than 1\%, the planned
caesarean probabilities by less than 0.005, and the Stage 5 parameters by less
than 3\%. Each round uses ten runs of 24 months, each with its own parameter
draw. The calibration converged in seven rounds.

\section{Results}

\begin{table}[htbp]
\centering
\caption{Estimated parameters.}
\label{tab:params}
\begin{tabular}{lrrr}
\toprule
 & Kisii & Makueni & Mombasa \\
\midrule
\p{PT\_scale} & 0.727 & 0.701 & 0.716 \\
\p{home\_noANC} & 0.1324 & 0.1076 & 0.0358 \\
\p{home\_lowrisk} & 0.0424 & 0.2241 & 0.1008 \\
\p{baseline\_l23\_preference} & 0.3214 & 0.3904 & 0.4563 \\
\p{p\_elec\_CS|highrisk} & 0.2517 & 0.2885 & 0.4173 \\
\p{p\_elec\_CS|preterm} (shared) & 0.1001 & 0.1001 & 0.1001 \\
\p{p\_aph} & 0.01533 & 0.02128 & 0.02254 \\
\p{p\_ruptured\_uterus} & 0.00180 & 0.00247 & 0.00266 \\
\p{p\_eclampsia} & 0.00955 & 0.01589 & 0.01171 \\
\p{p\_OL\_scale} & 1.242 & 1.995 & 2.312 \\
\p{pph\_incidence\_multiplier} & 0.766 & 0.883 & 0.792 \\
\p{sepsis\_incidence\_multiplier} & 0.0303 & 0.0414 & 0.0287 \\
\p{p\_cs\_capacity} [L2/3] & 0.0610 & 0.0661 & 0.0847 \\
\p{p\_cs\_capacity} [L4] & 0.1970 & 0.2883 & 0.3029 \\
\p{t\_l23\_l45\_preterm} (\%) & 71.0 & 73.1 & 73.8 \\
\p{t\_l23\_l45\_notsevere} (\%) & 55.0 & 47.8 & 54.5 \\
\p{t\_l23\_l45\_predicted} (\%) & 42.8 & 36.7 & 41.4 \\
\p{p\_mat\_death\_facility\_baseline} (pooled) & 0.000941 & 0.000941 & 0.000941 \\
\p{RR\_MM\_l23\_vs\_l4} (pooled) & 1.945 & 1.945 & 1.945 \\
\bottomrule
\end{tabular}
\end{table}

\begin{longtable}{lrrr}
\caption{Fit to targets (model / target).}\label{tab:fit}\\
\toprule
 & Kisii & Makueni & Mombasa \\
\midrule
\endfirsthead
\toprule
 & Kisii & Makueni & Mombasa \\
\midrule
\endhead
Final L2/3 share & 0.165 / 0.166 & 0.169 / 0.168 & 0.210 / 0.212 \\
Preterm among facility births & 0.133 / 0.134 & 0.134 / 0.134 & 0.135 / 0.134 \\
Planned CS at L4 & 0.080 / 0.080 & 0.112 / 0.111 & 0.125 / 0.125 \\
APH, L4 & 0.0146 / 0.0146 & 0.0204 / 0.0207 & 0.0220 / 0.0223 \\
Uterine rupture, L4 & 0.0018 / 0.0018 & 0.0023 / 0.0025 & 0.0025 / 0.0027 \\
Eclampsia, L4 & 0.0124 / 0.0125 & 0.0181 / 0.0176 & 0.0191 / 0.0190 \\
Obstructed labour, L4 & 0.058 / 0.058 & 0.082 / 0.082 & 0.089 / 0.088 \\
PPH, L4 & 0.049 / 0.049 & 0.070 / 0.070 & 0.075 / 0.075 \\
Sepsis, L4 & 0.0048 / 0.0047 & 0.0067 / 0.0067 & 0.0072 / 0.0072 \\
CS, L4 & 0.202 / 0.200 & 0.277 / 0.277 & 0.312 / 0.312 \\
CS, L2/3 & 0.040 / 0.040 & 0.066 / 0.065 & 0.086 / 0.085 \\
L2/3 preterm share & 0.060 / 0.059 & 0.063 / 0.062 & 0.064 / 0.064 \\
L2/3:L4 complication ratio & 0.391 / 0.393 & 0.571 / 0.573 & 0.512 / 0.506 \\
\midrule
Planned CS among preterm, pooled & \multicolumn{3}{c}{0.0728 / 0.0729} \\
\quad county values (checks) & 0.061 & 0.082 & 0.080 \\
\bottomrule
\end{longtable}

\begin{table}[htbp]
\centering
\caption{Maternal mortality, direct causes (deaths per 100,000 deliveries;
KHIS scaled by 0.618).}
\label{tab:mortality}
\begin{tabular}{llrrrc}
\toprule
County & Level & KHIS deaths & KHIS MMR (95\% CI) & Model & Within CI \\
\midrule
Pooled (fitted) & L2/3 & 7 & 7.1 & 7.4 & -- \\
Pooled (fitted) & L4 & 84 & 30.8 & 31.0 & -- \\
\midrule
Kisii & L2/3 & 2 & 6.1 (0.7--22.0) & 5.6 & yes \\
Kisii & L4 & 18 & 18.8 (11.1--29.7) & 24.5 & yes \\
Makueni & L2/3 & 4 & 16.6 (4.5--42.5) & 7.9 & yes \\
Makueni & L4 & 35 & 44.6 (31.1--62.0) & 34.4 & yes \\
Mombasa & L2/3 & 1 & 2.4 (0.1--13.5) & 8.6 & yes \\
Mombasa & L4 & 31 & 31.4 (21.3--44.5) & 35.5 & yes \\
\bottomrule
\end{tabular}
\end{table}

\section{Validation}

Validation uses 10 runs of 60 months, each with its own parameter draw and
simulation seeds not used in fitting. All 45 county targets and the three
pooled targets (planned caesarean section among preterm births, and L2/3 and
internal-L4 MMR) are reproduced within 10\%. The initial L2/3 share is implied
by the final share and $\tau$, so it is reported but not counted as a target.
Quantities not used in fitting serve as checks:
\begin{itemize}
\item County MMR lies within its exact Poisson 95\% interval in all six
county-level comparisons (Table~\ref{tab:mortality}); Mombasa's L2/3 interval
rests on one death.
\item Facility SMO is 0.031 / 0.045 / 0.047 against reference values of 0.036 /
0.040 / 0.044.
\item The share of L4/5 births arriving by transfer is 0.11 / 0.16 / 0.22,
against 0.122 in Waiswa et al.
\item PPH at L4 is 5--8\%, consistent with suspected PPH of 9\% (7--11\%) in
Kenyan referral hospitals.
\item The DHS L2/3 share among ANC mothers agrees in Kisii and Mombasa but not
in Makueni.
\item The KHIS L4 preterm rate ($\times 4.48$) is 0.088 / 0.160 / 0.181 against
0.145 / 0.147 / 0.151 in the model; Kisii's low recorded rate is the
under-recording noted for Stage 1.
\end{itemize}

\section{Sensitivity analyses}
\label{sec:sensitivity}

Each analysis repeats the full calibration and validation. Intervention effects are
compared with the same seeds (40 paired runs of 36 months) as relative changes
in expected MMR.

\begin{longtable}{>{\raggedright\arraybackslash}p{0.27\textwidth} >{\raggedright\arraybackslash}p{0.65\textwidth}}
\caption{Sensitivity analyses.}\label{tab:sensitivity}\\
\toprule
Analysis & Result \\
\midrule
\endfirsthead
\toprule
Analysis & Result \\
\midrule
\endhead
L2/3 monitoring without the 0.8 factor on \p{sen\_comp\_trad}, \p{spec\_comp\_trad} & $\tau$ unchanged; all targets met. \p{t\_l23\_l45\_predicted} rises to keep the Rescue share. The two values are kept; they set the comparator for the PULSE sensor intervention rather than any calibration target. \\
Transfer share fixed at 0.15 or 0.25, with transfers set from the full Rescue reason mix & L2/3 preterm share doubles (0.12 against KHIS 0.06) and the referral share of L4/5 births falls to 0.04--0.09; baseline complications and mortality barely change. At $\tau=0.25$, Transfer \& EMT effects at high implementation are similar in Kisii ($-5.0\%$ vs $-4.8\%$) and smaller in Makueni ($-3.9\%$ vs $-3.0\%$) and Mombasa ($-4.6\%$ vs $-2.9\%$); PROMPTS effects are larger where fewer women are transferred out of L2/3 (Makueni $-25.9\%$ vs $-29.4\%$); the combined package changes by at most 4 points. \\
$\theta=2.5$ & Parameters and county MMR change by less than 2\%; PROMPTS, Transfer \& EMT, and combined effects change by at most 2 percentage points. \\
County cause-specific complication targets & All targets met, but the PPH and sepsis multipliers vary 4- to 9-fold across counties. County L4 mortality fits KHIS slightly better (mean absolute error 16\% vs 22\%; the likelihood of the observed deaths is about 1.8 times higher), which is weak evidence given 18--35 deaths per county. Effects of interventions acting through specific complications shift in opposite directions in Kisii and Makueni, whose recorded cause mixes differ most from the pooled mix: PPH bundle $-10.3\%$ vs $-7.4\%$ in Kisii and $-10.0\%$ vs $-13.3\%$ in Makueni; PROMPTS $-27.2\%$ vs $-31.0\%$ and $-25.9\%$ vs $-19.3\%$; the combined package $-64.9\%$ vs $-59.9\%$ in Kisii. Mombasa changes by less than 1.2 points. \\
External referrals to L5 (illustrative) & L5 also delivers women referred from outside the county, who are not modeled; their number and acuity are not observed. For assumed external shares of L5 deliveries (0--30\%) and late-referral odds ratios (1--3), the reported L5 MMR is $(1-s)\times$ internal L5 MMR $+\ s\times$ proxy external MMR, the proxy being complication-bearing emergency transfers received at L5. One observed MMR cannot identify two unknowns, so the grid is shown rather than a fit. Against KHIS L5 MMR (direct causes) of 160 (125--202) in Kisii and 312 (259--373) in Mombasa, several combinations fit (Kisii 5--10\% external; Mombasa 10--30\%). In Makueni the modeled internal L5 MMR (177) already exceeds the observed 61 (38--93). \\
\bottomrule
\end{longtable}

\section{Applying the procedure to another county}

The procedure requires, for a new county, the DHS county estimates and KHIS
MOH 711 facility data listed in Table~\ref{tab:sources}, both of which exist for
all Kenyan counties, and a check of facility level coding. The remaining
inputs are literature values or are pooled. Three qualifications apply. First,
the transfer-reason mix, receiving level, and under-recording factor come from
a single county; a county with its own referral records should use them. Second, the
cause mix and the mortality level are pooled; adding counties strengthens the
pool, but a single new county is checked against, not fitted to, its own
mortality. Third, data quality varies by county (extract completeness, level
coding, small DHS samples) and should be reviewed before targets are built.

\section{Limitations}

\begin{enumerate}
\item The transfer share exceeds published primary-facility referral rates; it
is the value required to reproduce KHIS L2/3 composition and the Rescue reason
mix, rather than an observed rate.
\item The Rescue reason mix and receiving level come from one county.
\item \p{t\_l23\_l45\_severe}, \p{sen\_comp\_trad}, and \p{spec\_comp\_trad} have
no direct data.
\item A single under-recording factor is applied to all diagnoses and levels.
\item SMO matches its reference values on average over parameter draws but not
for individual draws with a low severity relative risk.
\item Makueni's DHS facility split conflicts with KHIS.
\item The direct-cause share of deaths is a global value.
\item \p{sepsis\_incidence\_multiplier} is 0.03--0.04: matching KHIS sepsis
($\times 4.48$) requires reducing the model's baseline sepsis risk about 30-fold,
which suggests that sepsis is under-recorded by more than the common factor.
\end{enumerate}

\end{document}
